\section{HW10: C Pointers + LDR/STR + Function Calls}

\subsection{C precedence (rows that matter)}
\begin{tabular}{@{}clll@{}}
\toprule
L & ops & assoc & note \\
\midrule
1 & \cee{() [] . -> ++ -- (post)} & L$\to$R & postfix inc/dec \\
2 & \cee{++ -- (pre)}, \cee{* \& (cast)} & R$\to$L & deref, addr-of \\
14 & \cee{= += ...} & R$\to$L & assignment \\
\bottomrule
\end{tabular}
\textbf{Rule:} postfix \cee{p++} (L1) yields OLD pointer; deref/prefix-\cee{++}/cast share L2 R$\to$L. Decode right-to-left.

\subsection{Idiom $\to$ ASM cookbook}
\begin{tabular}{@{}lll@{}}
\toprule
C & ASM & note \\
\midrule
\cee{v=*(p+k)}  & \asm{ldrb r0,[r1,\#k]}   & basic, p unchanged \\
\cee{v=*++p}    & \asm{ldrb r0,[r1,\#1]!}  & pre (p first) \\
\cee{v=*p++}    & \asm{ldrb r0,[r1],\#1}   & post (load first) \\
\cee{*(p+k)=v}  & \asm{strb r0,[r1,\#k]}   & basic store \\
\cee{*p++=v}    & \asm{strb r0,[r1],\#1}   & post store \\
\cee{++(*p)}    & ldr;add\#1;str           & READ-MOD-WRITE \\
\cee{(*p)++}    & ldr;mov;add\#1;str       & return old \\
\cee{++*p++}    & ldr post;add;str\asm{[\#-1]} & WB to OLD \\
\cee{*p++=*q++} & ldr post; str post       & two ptrs walk \\
\bottomrule
\end{tabular}
\textbf{Hook:} \texttt{!}$=$urgent update BEFORE; trailing \texttt{\#}$=$update AFTER.

\subsection{Type suffix table (load only --- stores zero/sign agnostic)}
\begin{tabular}{@{}llll@{}}
\toprule
sfx & C type & ext & sample \\
\midrule
(none) & \cee{int32\_t} & --- & \asm{ldr} \\
B  & \cee{uint8\_t}  & zero $\to$ 32 & 0x80$\to$0x0000\_0080 \\
SB & \cee{int8\_t}   & sign $\to$ 32 & 0x80$\to$0xFFFF\_FF80 \\
H  & \cee{uint16\_t} & zero $\to$ 32 & 0xB0A0$\to$0x0000\_B0A0 \\
SH & \cee{int16\_t}  & sign $\to$ 32 & 0xB0A0$\to$0xFFFF\_B0A0 \\
\bottomrule
\end{tabular}
Sign-extension test: MSB of the loaded chunk (bit 7 for SB, bit 15 for SH).
\textbf{STR has no S/U variant} (you store the bits you have).

\subsection{Pointer arithmetic scaling}
Increment of typed pointer scales by \cee{sizeof(*p)}: \cee{p8++}$=$+1B,
\cee{p16++}$=$+2B, \cee{p32++}$=$+4B. Register-offset shift exponent matches
log$_2$(elem): \asm{[Rb,Ri,LSL \#0/1/2/3]} for byte/hw/word/dword.
\textbf{Typed difference} \cee{q-p} returns \emph{elements}, not bytes
(divide-by-size baked in).

\subsection{HW10 P1 --- four-line memory \& pointer trace}
Array: \cee{uint8\_t my\_arr[16]} at \texttt{0x2000\_0100}, \cee{my\_arr[i]=2i}.
\cee{p8a=0x...100}, \cee{p8b=0x...10C}.
\begin{tabular}{@{}lllll@{}}
\toprule
ln & RHS & memory effect & p8a$\to$ & p8b$\to$ \\
\midrule
4 & \cee{*p8b++=*(p8a+2)}     & rd 0x102$=$0x04; \texttt{0x10C}$\leftarrow$0x04 & 0x100 & 0x10D \\
5 & \cee{*p8b++=++*(p8a+4)}   & 0x104:0x08$\to$0x09; \texttt{0x10D}$\leftarrow$0x09 & 0x100 & 0x10E \\
6 & \cee{*p8b++=++*p8a++}     & 0x100:0x00$\to$0x01; \texttt{0x10E}$\leftarrow$0x01 & 0x101 & 0x10F \\
7 & \cee{*p8b++=*p8a++}       & rd 0x101$=$0x02; \texttt{0x10F}$\leftarrow$0x02 & 0x102 & 0x110 \\
\bottomrule
\end{tabular}

\textbf{Bytes at 0x10C--0x10F:} \texttt{0x04,0x09,0x01,0x02}.
\textbf{Mem changes by RHS:} L4 none, L5 0x104(0x08$\to$0x09), L6 0x100(0x00$\to$0x01), L7 none.

\textbf{Line 6 dissection (\cee{++*p8a++}, the trickiest):}
postfix \cee{p8a++} (L1) gives OLD pointer 0x100; deref reads 0x00; prefix
\cee{++} adds 1, writes 0x01 BACK to 0x100. So memory at 0x100 changes; the
loaded value is 0x01; p8a now 0x101.

\subsection{HW10 P1.3 --- ASM (leaf, \cee{my\_fn(uint8\_t*, uint8\_t*)})}
\begin{lstlisting}
my_fn:                       @ R0=p8a, R1=p8b (caller-saved)
  @ L4: *p8b++ = *(p8a+2);
    ldrb r2, [r0, #2]          @ basic offset, p8a stays
    strb r2, [r1], #1          @ post-idx: *p8b=r2; p8b++
  @ L5: *p8b++ = ++*(p8a+4);
    ldrb r2, [r0, #4]
    add  r2, r2, #1            @ prefix-inc the value
    strb r2, [r0, #4]          @ WRITEBACK -- don't forget!
    strb r2, [r1], #1
  @ L6: *p8b++ = ++*p8a++;
    ldrb r2, [r0], #1          @ post-idx: r2=*OLD p8a; R0+=1
    add  r2, r2, #1
    strb r2, [r0, #-1]         @ writeback to OLD addr (R0-1)
    strb r2, [r1], #1
  @ L7: *p8b++ = *p8a++;
    ldrb r2, [r0], #1
    strb r2, [r1], #1
    bx   lr                     @ leaf: no PUSH/POP, no LR save
\end{lstlisting}

\textbf{Three patterns to memorize:}
(1) \cee{++(*(p+k))} $=$ ldrb / add\#1 / strb-to-same-offset / strb-to-dest;
(2) \cee{++*p++}     $=$ ldrb post-idx / add / strb \asm{[\#-1]} / strb post-idx;
(3) \cee{*p++}       $=$ ldrb post-idx alone (no writeback).

\subsection{AAPCS: who-saves-what (the only table you need)}
\begin{tabular}{@{}lll@{}}
\toprule
reg & role & preserved? \\
\midrule
R0--R3 & args + return value & no (caller-saved, scratch) \\
R4--R11& callee scratch       & \textbf{yes} (PUSH if used) \\
R12 (IP)& scratch              & no \\
R13 (SP)& stack pointer        & yes (8-byte aligned at boundary) \\
R14 (LR)& return addr          & save if you BL (stem) \\
R15 (PC)& program counter      & --- \\
\bottomrule
\end{tabular}

\textbf{Caller widens narrow types} (\cee{int8\_t/int16\_t}) to 32b before
placing in R0--R3: signed args sign-ext, unsigned zero-ext. 64-bit args use
register pair (lo, hi), even-aligned (R0:R1 or R2:R3).

\subsection{Leaf vs Stem function templates}
\textbf{Leaf} (no BL inside, $\le$4 scratch regs needed): no push, no pop.
\begin{lstlisting}
leaf_fn:                  @ R0..R3 in, R0 out
    @ work in R0--R3, R12
    bx   lr
\end{lstlisting}

\textbf{Stem} (calls others or needs R4--R11): push LR + any callee-saved used.
Stash-args pattern: copy R0--R3 into R4--R7 \emph{before} BL clobbers them.
\begin{lstlisting}
stem_fn:
    push {r4-r7, lr}        @ save callee + return addr
    mov  r4, r0             @ stash args (R0..R3 die on next BL)
    mov  r5, r1
    mov  r6, r2
    mov  r7, r3
    @ ... bl other_fn; arg in R0; result in R0 ...
    @ recover args from R4--R7 between calls
    pop  {r4-r7, pc}        @ pop PC = return; replaces BX LR
\end{lstlisting}
Push/pop list \emph{must match}; pushing odd \# of regs $\to$ pad for 8B align
(usually push LR alongside even count). \asm{POP \{...,pc\}} returns directly.

\subsection{HW10 P2 --- four-task LDRD/SB/SH/H trace}
Setup: \cee{ipt[i]=i<<4} $\to$ bytes 0x00,0x10,0x20,...,0xF0. Word at ipt+0
$=$ 0x30201000 (little-endian: LSB at lowest addr). Halfwords at
ipt+0/+2/.../+14 $=$ 0x1000,0x3020,0x5040,0x7060,0x9080,0xB0A0,0xD0C0,0xF0E0.
Entry R0$=$0x2000\_1010 (ipt), R1$=$0x2000\_1040 (opt).

\textbf{T1: LDRD pre-idx (8B$\to$2 regs).}
\asm{ldrd r2,r3,[r0,\#8]!} : R0$+$=8 first, then load lo word at \texttt{[R0]}
$=$ 0xB0A09080, hi word at \texttt{[R0+4]} $=$ 0xF0E0D0C0.
\asm{strd r2,r3,[r1,\#0]} writes opt[0],opt[1].
\textbf{Print:} \texttt{Out1=0xB0A09080, Out2=0xF0E0D0C0}.

\textbf{T2: LDRSB pre-idx, sign-extend.}
\asm{ldrsb r2,[r0,\#4]!}: R0$\to$ipt+4, byte 0x40 (bit7$=$0) $\to$ R2$=$0x0000\_0040.
\asm{ldrsb r3,[r0,\#4]!}: R0$\to$ipt+8, byte 0x80 (bit7$=$1) $\to$ R3$=$0xFFFF\_FF80.
After: R0$=$\textbf{0x2000\_1018}, R1$=$0x2000\_1040 (STRD has no \texttt{!}).
\textbf{Print:} \texttt{Out1=0x40, Out2=0xFFFFFF80}.

\textbf{T3: LDRSH post-idx, odd stride 10.}
\asm{ldrsh r2,[r0],\#10}: load hw at ipt+0$=$0x1000 (bit15$=$0) $\to$
R2$=$0x0000\_1000; THEN R0$+$=10.
\asm{ldrsh r3,[r0],\#4}: hw at ipt+10$=$0xB0A0 (bit15$=$1) $\to$
R3$=$0xFFFF\_B0A0; R0$+$=4.
After: R0$=$\textbf{0x2000\_101E} (unaligned!).
\textbf{Print:} \texttt{Out1=0x1000, Out2=0xFFFFB0A0}.

\textbf{T4: LDRH register-offset, zero-ext.} Called \cee{task4(ipt,opt,1)}, R2$=$1.
\asm{ldrh r3,[r0,r2,LSL \#1]}: EA$=$ipt$+$2, hw 0x3020 $\to$ R3$=$0x0000\_3020;
\asm{str r3,[r1,\#0]}; \asm{add r2,\#1} (R2$=$2); reload at ipt$+$4 $=$
0x5040 $\to$ R3$=$0x0000\_5040; \asm{str r3,[r1,\#4]}; \asm{add r2,\#1} (R2$=$3).
\textbf{Print:} \texttt{Out1=0x3020, Out2=0x5040}. \textbf{R2$=$0x3, R3$=$0x5040.}

\subsection{Sign-extend pivots: LOAD vs CALL}
\textbf{LOAD:} LDRSB/LDRSH test bit 7/15 of loaded data, widen to 32b. LDRB/LDRH zero-pad.
\textbf{CALL:} caller widens narrow C arg per type signedness BEFORE \asm{BL}. Two independent extensions.

\subsection{HW10 P2 --- ASM$\to$C reverse}
\textbf{T1} 8B copy: \cee{*((int64\_t*)opt)=*((int64\_t*)ipt+1)} (or word-by-word with \cee{(int32\_t*)}).
\textbf{T2} pre-inc \cee{int8\_t*}: \cee{p8+=4; opt[0]=*p8; p8+=4; opt[1]=*p8;}.
\textbf{T3} post-inc \cee{int16\_t*}: \cee{opt[0]=*p16; opt[1]=*(p16+5);} (\#10$=$5 elem).
\textbf{T4} idx \cee{uint16\_t*}: \cee{int i=1; opt[0]=*(p16u+i); i++; opt[1]=*(p16u+i); i++;}.

\subsection{Array-loop $+$ call pattern (compute index $\to$ ldr/bl/str)}
Idiom: \cee{for(i;..) arr[i]=fn(arr[i],...)} compiles to:
\begin{lstlisting}
    @ R5=&arr (base, callee-saved), R8=i (callee-saved)
loop:
    ldr  r0, [r5, r8, lsl #2]   @ R0 = arr[i]   (32b elements)
    bl   fn                      @ R0 = result; clobbers R0--R3,R12
    str  r0, [r5, r8, lsl #2]   @ arr[i] = R0
    add  r8, r8, #1
    cmp  r8, r9                  @ R9 = N (callee-saved)
    blt  loop
\end{lstlisting}
\textbf{Why callee-saved (R4--R11) for base/index/N:} they survive across
\asm{BL fn}. Putting them in R0--R3/R12 would corrupt on every call.

\subsection{Fast-failure modes}
\begin{itemize}
\item Forgetting STRB writeback after \cee{++(*p)}: needs \asm{ldr;add;str} then dest store.
\item \asm{[r0,\#4]!} (pre, urgent BEFORE) vs \asm{[r0],\#4} (post, AFTER).
\item LDRB on \cee{int8\_t} $\to$ wrong zero-pad; use LDRSB.
\item STRD/LDRD w/o \texttt{!} doesn't update Ra (T1 R1 unchanged; T2 R0 advances).
\item Reg-offset shift: LSL \#0/1/2/3 for byte/hw/word/dword.
\item No \asm{push \{lr\}} in stem $\to$ \asm{BL} clobbers LR, return broken.
\end{itemize}

\subsection{One-line summary table}
\begin{tabular}{@{}lll@{}}
\toprule
goal & instruction & C \\
\midrule
\cee{*(p+k)}      & \asm{ldr\{B/H\}\{S\} Rd,[Ra,\#os]} & basic, p stays \\
\cee{*++p}        & \asm{...,[Ra,\#os]!}                & pre, p$+$=os first \\
\cee{*p++}        & \asm{...,[Ra],\#os}                 & post, load then p$+$=os \\
\cee{*(p+i)} typed& \asm{...,[Ra,Ri,LSL \#k]}           & k$=$log$_2$elem \\
8B copy           & \asm{ldrd}/\asm{strd}               & 64b, two regs \\
sign-ext load     & \asm{ldrsb}/\asm{ldrsh}             & test bit 7/15 \\
leaf return       & \asm{bx lr}                          & no push \\
stem return       & \asm{pop \{...,pc\}}                & pushed LR \\
\bottomrule
\end{tabular}
